% =============================================================================
\section{更广泛的背景：相关工作}
\label{sec:related}
% =============================================================================

\paragraph{综述。}
Moslem 和 Kelleher~\cite{moslem2026routingsurvey}综述了动态模型路由与级联，覆盖查询难度、人类偏好、聚类、不确定性、RL 与多模态路由；但他们关于何时/什么/如何的分类，并未涉及资源池架构或工作负载交互。Pan 与 Li~\cite{pan2025survey}覆盖了从算子优化到无服务器部署的完整推理栈，但把路由视为集群级负载均衡。Xia~\etal~\cite{xia2025taming}则从实例级和集群级两个层面综述高效 LLM 推理服务。这些综述都没有形式化 WRP 框架所关注的跨维交互。

\paragraph{路由系统（不含我们的项目）。}
RouteLLM~\cite{ong2025routellm}与 Hybrid LLM~\cite{ding2024hybrid}使用偏好数据在强/弱模型之间做路由。FrugalGPT~\cite{chen2023frugalgpt}与 AutoMix~\cite{madaan2024automix}使用级联。Router-R1~\cite{routerr1_2025}与 R2-Router~\cite{r2router2026}使用基于 RL 的选择。MixLLM~\cite{lu2025mixllm}在 NAACL 2025 上使用上下文 bandit。xRouter~\cite{xrouter2025}则在 20+ LLM 工具上训练成本感知 RL 路由器，并显式使用成本-性能奖励。

\paragraph{智能体服务与失败分析。}
AGSERVE~\cite{agserve2025}为智能体工作负载提供会话感知的 KV-cache 管理与模型级联（NeurIPS 2025）。Helium~\cite{helium2026}把智能体工作流当作查询计划来处理，并引入主动缓存和缓存感知调度。SUTRADHARA~\cite{sutradhara2026}围绕工具型智能体，对编排器与服务引擎做联合设计。Continuum~\cite{continuum2025}在 SWE-Bench 和 BFCL 上为多轮智能体调度引入 KV-cache TTL pinning。Concur~\cite{concur2026}通过基于拥塞的并发调节，把智能体级接纳控制引入系统（ICML 2026）。与此同时，越来越多的经验研究开始分类智能体故障模式~\cite{agenticfaults2026}（从 13,602 个问题中归纳出 37 种故障类型）、模型特定失败特征~\cite{erroratlas2026}（83 个模型、35 个数据集），以及多轮模型切换退化~\cite{handoffdrift2026}（一次切换即可带来 $\pm$4--13 个百分点波动）；但这些工作都没有把失败信号反馈回路由决策之中。

\paragraph{集群配置（不含我们的项目）。}
M\'{e}lange~\cite{griggs2024melange}优化异构 GPU 分配。SageServe~\cite{jia2025sageserve}则在微软规模上把流量预测整合进来。

\paragraph{解耦推理。}
Splitwise~\cite{patel2024splitwise}、DistServe~\cite{zhong2024distserve}与 Sarathi-Serve~\cite{agrawal2024sarathi}把 prefill 与 decode 解耦。Mooncake~\cite{qin2025mooncake}提供以 KV-cache 为中心的解耦。NVIDIA Dynamo~\cite{nvidia2026dynamo}实现了动态多节点调度。

\paragraph{能效（不含我们的项目）。}
SweetSpot~\cite{patel2026sweetspot}提供分析式能耗预测。TokenPowerBench~\cite{niu2025tokenpowerbench}与 GreenServ~\cite{greenserv2026}则对推理能耗进行基准测试并提出优化方法。
